% STAGE08-CHAPTER: V1-C07
% CHAPTER-SCHEMA-COVERAGE: CH-01,CH-02,CH-03,CH-04,CH-05,CH-06,CH-07,CH-08,CH-09,CH-10,CH-11,CH-12,CH-13,CH-14,CH-15,CH-16,CH-17,CH-18,CH-19,CH-20,CH-21,CH-22,CH-23,CH-24,CH-25,CH-26,CH-27,CH-28,CH-29,CH-30,CH-31,CH-32,CH-33,CH-34,CH-35,CH-36,CH-37
% CHAPTER-SCHEMA-EVIDENCE: CH-01=\label{struct:V1-C07-CH01}; CH-02=\label{struct:V1-C07-CH02}; CH-03=\label{struct:V1-C07-CH03}; CH-04=\label{struct:V1-C07-CH04}; CH-05=\label{struct:V1-C07-CH05}; CH-06=\label{struct:V1-C07-CH06}; CH-07=\label{struct:V1-C07-CH07}; CH-08=\label{struct:V1-C07-CH08}; CH-09=\label{struct:V1-C07-CH09}; CH-10=\label{struct:V1-C07-CH10}; CH-11=\label{struct:V1-C07-CH11}; CH-12=\label{struct:V1-C07-CH12}; CH-13=\label{struct:V1-C07-CH13}; CH-14=\label{struct:V1-C07-CH14}; CH-15=\label{struct:V1-C07-CH15}; CH-16=\label{struct:V1-C07-CH16}; CH-17=\label{struct:V1-C07-CH17}; CH-18=\label{struct:V1-C07-CH18}; CH-19=\label{struct:V1-C07-CH19}; CH-20=\label{struct:V1-C07-CH20}; CH-21=\label{struct:V1-C07-CH21}; CH-22=\label{struct:V1-C07-CH22}; CH-23=\label{struct:V1-C07-CH23}; CH-24=\label{struct:V1-C07-CH24}; CH-25=\label{struct:V1-C07-CH25}; CH-26=\label{struct:V1-C07-CH26}; CH-27=\label{struct:V1-C07-CH27}; CH-28=\label{struct:V1-C07-CH28}; CH-29=\label{struct:V1-C07-CH29}; CH-30=\label{struct:V1-C07-CH30}; CH-31=\label{struct:V1-C07-CH31}; CH-32=\label{struct:V1-C07-CH32}; CH-33=\label{struct:V1-C07-CH33}; CH-34=\label{struct:V1-C07-CH34}; CH-35=\label{struct:V1-C07-CH35}; CH-36=\label{struct:V1-C07-CH36}; CH-37=\label{struct:V1-C07-CH37}

% FIGURE-KNOWLEDGE-MAP: fig:V1-C07-convex-set -> PRE-OP-006
\chapter{凸优化与拉格朗日对偶}\label{chap:V1-C07}
\SLChapterOpening{从可行域直达凸性、对偶与KKT证书}{约束优化何时能由局部条件确认全局最优}{能够检查五类KKT条件并解释强对偶前提}{\cref{chap:V1-C02,chap:V1-C03}的正定性、梯度和Hessian}{资格条件或乘子符号不清楚时返回原问题与拉格朗日节}
\index{凸优化}

\index{拉格朗日对偶}


\phantomsection\label{struct:V1-C07-CH01}
\chaptergoals{
\begin{itemize}[leftmargin=2em]
  \item 严格区分可行性、局部最优、全局最优、凸性与严格凸性；
  \item 从原问题逐步构造拉格朗日函数、对偶函数和对偶问题，证明弱对偶并说明强对偶条件；
  \item 能够写出并检验KKT条件，解释互补松弛对应的活动约束。
\end{itemize}}

\phantomsection\label{struct:V1-C07-CH02}
\begin{prerequisitebox}
需要向量、内积、正定矩阵、梯度、Hessian和方向导数。有限加权Jensen不等式已在\cref{pre:V1-C06-finite-jensen}随用随教；本章从凸性定义重新建立并证明它。概率模型中的最大似然或最大后验问题可通过取负号改写为本章的最小化形式。
\end{prerequisitebox}

\phantomsection\label{struct:V1-C07-CH03}
\begin{precheckbox}
\begin{enumerate}[leftmargin=2.2em]
  \item 能否由 $a^{\mathsf T}x\le b$ 判断半空间是否凸？
  \item 能否计算二次函数的梯度和Hessian，并用特征值判断半正定？
  \item 能否说明“约束满足”与“目标值最小”是两个不同判断？
\end{enumerate}
未通过时回看\cref{sec:V1-C02-S06,sec:V1-C03-S02,sec:V1-C03-S03}。
\end{precheckbox}

\phantomsection\label{struct:V1-C07-CH04}
\begin{dependencybox}
可行域 $+$ 目标函数 $\to$ 优化问题；
凸集 $+$ 凸函数 $\to$ 局部最优等于全局最优；
约束 $\to$ 拉格朗日函数 $\to$ 对偶函数 $\to$ 弱/强对偶 $\to$ KKT与互补松弛。
\end{dependencybox}

\phantomsection\label{struct:V1-C07-CH05}
% STAGE19-SYMBOL-INDEX-BEGIN: V1-C07
\index[symbols]{x-minus-in-minus-minus-mathbbr-minus-d--s37a7f1741e55@$x\in\mathbb R^d$：待优化变量}
\index[symbols]{f-x--sfcd22889b6e7@$f(x)$：目标函数}
\index[symbols]{c-i-x-h-j-x--s52e0dbf7f88a@$c_i(x),h_j(x)$：不等式约束与等式约束函数}
\index[symbols]{minus-alpha-minus-i-minus-beta-minus-j--sf31d3cc93a19@$\alpha_i,\beta_j$：拉格朗日乘子}
\index[symbols]{l-x-minus-alpha-minus-minus-beta-minus--sef75028887a8@$L(x,\alpha,\beta)$：拉格朗日函数}
\index[symbols]{p-u002a-d-u002a--sec83724304c8@$p^*,d^*$：原问题与对偶问题最优值}
% STAGE19-SYMBOL-INDEX-END: V1-C07
\begin{symboltablebox}
\begin{tabularx}{\linewidth}{@{}l l X@{}}
\toprule 符号 & 取值 & 含义\\ \midrule
$x\in\mathbb R^d$ & 决策向量 & 待优化变量\\
$f(x)$ & $\mathbb R$ & 目标函数\\
$c_i(x),h_j(x)$ & $\mathbb R$ & 不等式约束与等式约束函数\\
$\alpha_i,\beta_j$ & $\alpha_i\ge0,\beta_j\in\mathbb R$ & 拉格朗日乘子\\
$L(x,\alpha,\beta)$ & $\mathbb R$ & 拉格朗日函数\\
$p^*,d^*$ & 扩展实数 & 原问题与对偶问题最优值\\
\bottomrule
\end{tabularx}
\end{symboltablebox}

% KNOWLEDGE-PLAN: KN-V1-C07-CONCEPT-001 L1
\paragraph{读前自检：优化问题的基本形式。}约束优化标准形以决策$x$、目标$f(x)$及$c_i(x)\le0,h_j(x)=0$为对象；请把“最大化$g$、约束$a(x)\ge b$”分别改写为最小化$-g$和$b-a(x)\le0$，核对不等号方向统一。\par

\SLDirectSection{优化问题、目标函数与约束}{sec:V1-C07-S01}
\phantomsection\label{struct:V1-C07-CH06}
\knowledgeanchor{PRE-OP-001}{优化问题的基本形式}
% KNOWLEDGE-PLAN: KN-V1-C07-FORMULA-001 L1
\paragraph{读前自检：目标函数。}凸优化与拉格朗日对偶中的目标函数把问题限定在“目标函数把每个可行决策映射为可比较的标量”；请把“目标函数把每个可行决策映射为可比较的标量”中的已声明对象代回凸优化与拉格朗日对偶中的目标函数的定义，检查是否得到统计学习中常由经验损失、正则项或负对数似然组成。\par

\begin{definition}[约束优化问题]\label{struct:V1-C07-CH07}
给定目标函数、决策变量及一组等式和不等式约束，在全部满足约束的决策中寻找目标值最小者，称为约束最小化问题；若目标与不等式函数凸、等式函数仿射，则称为凸优化问题。
\end{definition}
本章采用约束最小化标准形式
\[
\min_{x\in\mathbb R^d} f(x)
\quad\text{s.t.}\quad c_i(x)\le0\ (i=1,\ldots,m),\quad h_j(x)=0\ (j=1,\ldots,r).
\]
最大化$g$可改写为最小化$-g$，而 $a(x)\ge b$ 可改写为 $b-a(x)\le0$。符号方向必须在构造乘子前统一。
\knowledgeanchor{PRE-OP-002}{目标函数}
目标函数把每个可行决策映射为可比较的标量。统计学习中常由经验损失、正则项或负对数似然组成。
\knowledgeanchor{PRE-OP-003}{等式约束与不等式约束}
等式约束把解限制在低维集合上；不等式约束定义带边界的可行区域。约束函数的尺度不改变可行集，却会改变对应乘子的尺度。

% KNOWLEDGE-PLAN: KN-V1-C07-CONCEPT-003 L1
\paragraph{读前自检：可行域。}在“任何不满足全部约束的点即使目标更小，也不能与可行最优点竞争”成立时处理凸优化与拉格朗日对偶中的可行域：请代入正文给出的取等边界检验$\mathcal C=\{x:c_i(x)\le0,\ h_j(x)=0\}$，并以不等号方向、边界取等和定义域彼此相容判定结果。\par
% KNOWLEDGE-PLAN: KN-V1-C07-CONCEPT-004 L1
\paragraph{读前自检：局部最优与全局最优。}局部最优$x^*$只要求某邻域内$f(x^*)\le f(x)$，全局最优则要求该式对整个可行域成立；请用同一$x^*$逐项比较两个量词范围，并核对全局最优必为局部最优、反向一般不成立。\par

\phantomsection\label{struct:V1-C07-CH08}
\textbf{模型。}模型由决策变量、目标函数和约束族组成；同一预测任务选择不同正则项，会得到不同的优化模型。
\phantomsection\label{struct:V1-C07-CH09}
\textbf{学习策略。}先统一最小化与“$\le0$”记号，再检查可行域非空、目标下界和约束资格，最后选择一阶或二阶算法。
\SLReviewSection{可行域与局部/全局最优}{sec:V1-C07-S02}
\knowledgeanchor{PRE-OP-004}{可行域}
可行域 $\mathcal C=\{x:c_i(x)\le0,\ h_j(x)=0\}$；任何不满足全部约束的点即使目标更小，也不能与可行最优点竞争。
\knowledgeanchor{PRE-OP-005}{局部最优与全局最优}
可行点$x^*$是局部最优，若存在半径$\delta>0$，对 $\mathcal C\cap B(x^*,\delta)$ 中所有点有 $f(x^*)\le f(x)$；若该关系对全部可行点成立，则为全局最优。


% KNOWLEDGE-PLAN: KN-V1-C07-CONCEPT-005 L1
\paragraph{读前自检：凸集。}围绕凸优化与拉格朗日对偶中的凸集，先采用条件“集合$C$凸是指任意$x,y\in C$与$t\in[0,1]$满足 $tx+(1-t)y\in C$”，再对$C$完成一项求导、代入或边界比较；最后验证凸集的任意交仍凸。\par
% KNOWLEDGE-PLAN: KN-V1-C07-PREREQUISITE-007 L1
\paragraph{读前自检：凸函数。}凸优化与拉格朗日对偶中的凸函数中的“函数$f:C\to\mathbb R$在凸域上凸”决定可执行范围；请对$f:C\to\mathbb R$完成一项求导、代入或边界比较，结果须满足结果的形状、符号和取等条件与定义一致。\par
% KNOWLEDGE-PLAN: KN-V1-C07-PREREQUISITE-008 L1
\paragraph{读前自检：严格凸函数。}把严格凸函数放在“函数$f:C\to\mathbb R$在凸域上凸”下考察：对$f:C\to\mathbb R$完成一项求导、代入或边界比较，并检查结果的形状、符号和取等条件与定义一致。\par

\SLTeachSection{凸集与凸函数}{sec:V1-C07-S03}
\knowledgeanchor{PRE-OP-006}{凸集}
集合$C$凸是指任意$x,y\in C$与$t\in[0,1]$满足 $tx+(1-t)y\in C$。仿射子空间、半空间、范数球均为凸集；凸集的任意交仍凸。
\knowledgeanchor{PRE-OP-007}{凸函数}
函数$f:C\to\mathbb R$在凸域上凸，若
\[
f(tx+(1-t)y)\le tf(x)+(1-t)f(y).
\]
\knowledgeanchor{PRE-OP-008}{严格凸函数}
若对$x\ne y$和$t\in(0,1)$上式严格小于，则$f$严格凸。严格凸目标在凸可行域上至多有一个全局最优点，但最优点存在仍需另行保证。

\phantomsection\label{struct:V1-C07-CH16}
% CORE-DERIVATION-PLAN: KN-V1-C07-DERIVATION-001
\SLKnowledgeBridge{逐步说明为何任意对偶可行值都不超过任意原可行值。}{任意原可行点 $x$ 与乘子 $\alpha\ge0,\beta$。}{不等式采用 $c_i(x)\le0$ 约定；等式约束满足 $h_j(x)=0$。}{先由可行性写 $L(x,\alpha,\beta)=f(x)+\sum_i\alpha_i c_i(x)+\sum_j\beta_jh_j(x)\le f(x)$。}

\begin{derivationgoalbox}
由定义逐步推出弱对偶，明确每个不等号的条件。
\end{derivationgoalbox}
\begin{derivationprepbox}
固定任意原可行$x$与任意$\alpha\ge0,\beta$；对偶函数定义为拉格朗日函数对$x$的下确界。
\end{derivationprepbox}
\textbf{推导第1步：}下确界不大于任一函数值，故 $g(\alpha,\beta)=\inf_zL(z,\alpha,\beta)\le L(x,\alpha,\beta)$。

\textbf{推导第2步：}原可行性给出$c_i(x)\le0,h_j(x)=0$，对偶可行性给出$\alpha_i\ge0$，所以$L(x,\alpha,\beta)\le f(x)$。

\textbf{推导第3步：}先对全部对偶可行乘子取上确界，再对全部原可行$x$取下确界，保持方向得到 $d^*\le p^*$。


% KNOWLEDGE-PLAN: KN-V1-C07-BOUNDARY-001 L1
\paragraph{读前自检：KKT条件。}KKT条件包含原可行、对偶可行、驻点与互补松弛；请对一个候选$(x^*,\alpha^*,\beta^*)$逐项代入四组关系，核对凸目标/不等式与仿射等式下满足这些条件即可推出全局最优。\par
% KNOWLEDGE-PLAN: KN-V1-C07-THEOREM-001 L1
\paragraph{读前自检：定理C 3。}定理C.3区分KKT的充分性、必要性与Slater作用；请核对“凸目标和不等式、仿射等式”使KKT点充分为全局最优，而从局部最优反推乘子还需约束资格，不能把两方向互换。\par

\SLAdvancedSection{KKT条件与互补松弛}{sec:V1-C07-S09}
\knowledgeanchor{PRE-OP-020}{KKT条件}
\knowledgeanchor{SLM-C-03-002}{定理C.3}
先分清三个方向。
\begin{enumerate}[leftmargin=2em]
  \item \textbf{充分性：}若目标与不等式约束凸、等式约束仿射，则任何原--对偶可行且满足驻点与互补松弛的点都是全局最优；这一方向不需要Slater条件。
  \item \textbf{必要性：}从“局部最优”推出乘子和KKT条件，需要线性无关、MFCQ、Slater等适用的约束资格，不能只凭凸性。
  \item \textbf{Slater的作用：}对凸问题，严格可行性保证强对偶；在本章的可微有限维设定及最优值可达到等常见正则前提下，它也保证最优点具有KKT乘子。
\end{enumerate}
因此，对可微约束问题，若局部最优点满足相应约束资格，则存在乘子使以下条件成立：
% KNOWLEDGE-PLAN: KN-V1-C07-BOUNDARY-002 L1
\paragraph{读前自检：互补松弛条件。}互补松弛要求$\alpha_i^*c_i(x^*)=0$且$\alpha_i^*\ge0,c_i(x^*)\le0$；请分别代入$c_i(x^*)<0$与$\alpha_i^*>0$，核对前者推出$\alpha_i^*=0$、后者推出约束取等。\par

\begin{align*}
\nabla_xL(x^*,\alpha^*,\beta^*)&=0 &&\text{驻点条件},\\
c_i(x^*)&\le0 &&\text{原可行性},\\
h_j(x^*)&=0 &&\text{原可行性},\\
\alpha_i^*&\ge0 &&\text{对偶可行性},\\
\alpha_i^*c_i(x^*)&=0 &&\text{互补松弛}.
\end{align*}
在凸性与仿射等式前提下，这些条件本身已经足以证明全局最优；Slater用于反方向的必要性与强对偶，不是该充分性证明的附加前提。
\knowledgeanchor{PRE-OP-021}{互补松弛条件}
若$c_i(x^*)<0$，则必须$\alpha_i^*=0$；若$\alpha_i^*>0$，则约束必须活动，即$c_i(x^*)=0$。乘子为零时约束可以活动也可以非活动。

% DERIVATION-ID: V1-C07-kkt-sufficiency-zero-gap
% CORE-DERIVATION-PLAN: KN-V1-C07-DERIVATION-002
\SLKnowledgeBridge{把驻点、可行性与互补松弛连成全局最优性证明。}{凸目标、凸不等式、仿射等式及候选原--对偶变量 $(x^*,\alpha^*,\beta^*)$。}{$x^*$ 原可行，$(\alpha^*,\beta^*)$ 对偶可行，并满足驻点与互补松弛。}{先由凸性和驻点证明 $x^*$ 是 $L(x,\alpha^*,\beta^*)$ 关于 $x$ 的全局极小点。}

\begin{derivationgoalbox}
证明在目标与不等式约束凸、等式约束仿射时，一组满足KKT条件的原--对偶变量给出全局最优解与零对偶间隙。
\end{derivationgoalbox}
\begin{derivationprepbox}
采用$c_i(x)\le0$、$\alpha_i\ge0$的符号约定；$x^*$原可行，$(\alpha^*,\beta^*)$对偶可行，并满足驻点与互补松弛。固定乘子后，$L(x,\alpha^*,\beta^*)$关于$x$为可微凸函数。
\end{derivationprepbox}
\textbf{推导第1步：由驻点得到拉格朗日函数的全局极小。}凸可微函数的一阶条件给出
\[
\nabla_xL(x^*,\alpha^*,\beta^*)=0
\quad\Longrightarrow\quad
x^*\in\arg\min_xL(x,\alpha^*,\beta^*).
\]
因此$g(\alpha^*,\beta^*)=L(x^*,\alpha^*,\beta^*)$。

\textbf{推导第2步：用可行性与互补松弛消去约束项。}等式可行性使$\beta_j^*h_j(x^*)=0$，互补松弛使$\alpha_i^*c_i(x^*)=0$，于是
\[
L(x^*,\alpha^*,\beta^*)=f(x^*).
\]

\textbf{推导第3步：夹逼原--对偶目标。}弱对偶对任意原、对偶可行点给出
\[
g(\alpha^*,\beta^*)\le d^*\le p^*\le f(x^*).
\]
前两步已给出左右端相等，故链中各量相等；$x^*$与$(\alpha^*,\beta^*)$分别全局最优，且$p^*-d^*=0$。

\phantomsection\label{struct:V1-C07-CH14}
\begin{corollary}[零间隙证书]
若凸问题的一组原--对偶变量满足KKT条件，则其原目标等于对偶目标，因而二者均最优。
\end{corollary}
% KNOWLEDGE-PLAN: KN-V1-C07-PROOF-001 L1
\paragraph{读前自检：证明：零间隙证书。}零间隙证书已由拉格朗日函数极小性、约束项消去和弱对偶夹逼得到$g(\alpha^*,\beta^*)=f(x^*)$；请写出夹逼链并核对等号同时证明原--对偶最优与零间隙。\par

\begin{proof}
前述三级推导分别建立拉格朗日函数极小性、约束项消去和弱对偶夹逼，已经得到$g(\alpha^*,\beta^*)=f(x^*)$，所以零间隙与原--对偶最优性成立。
\end{proof}

\phantomsection\label{struct:V1-C07-CH20}
\begin{example}[带非负约束的二次问题]\label{exm:V1-C07-kkt-nonnegative}
分别求$\min_x(x-2)^2$与$\min_x(x+1)^2$，共同约束为$x\ge0$，并比较非活动与活动约束的KKT乘子。
\end{example}
\SLExampleSolutionHeading{exm:V1-C07-kkt-nonnegative}
\begin{solution}
\SLReadTranslation{对两个具有共同非负约束的一维凸二次问题分别求KKT解，并比较约束非活动与活动时的乘子。}
\SolGiven 可行域为$[0,\infty)$，标准不等式写作$-x\le0$；两个目标均凸且存在严格可行点，KKT条件充分必要。
\SLMethodTrigger{分别写出拉格朗日函数和KKT四类条件，按互补松弛求解；再与无约束极小点及边界单调性比较。}
\SolPlan 对$f_1,f_2$各解一次驻点、原可行性、对偶可行性和互补条件，最后直接检查无约束候选是否落在$[0,\infty)$内。
\SolDerive
可行域为$\mathcal C=[0,\infty)$，标准约束$c(x)=-x\le0$满足Slater条件，例如$x=1$严格可行；两个目标均凸，因此KKT条件充分且必要。对$f_1(x)=(x-2)^2$，
\[
\begin{aligned}
L_1(x,\alpha)&=(x-2)^2-\alpha x,\qquad \alpha\ge0,\\
2(x-2)-\alpha&=0,\quad x\ge0,\quad \alpha x=0,\\
(x_1^\star,\alpha_1^\star)&=(2,0).
\end{aligned}
\]
对$f_2(x)=(x+1)^2$，
\[
\begin{aligned}
L_2(x,\alpha)&=(x+1)^2-\alpha x,\\
2(x+1)-\alpha&=0,\quad x\ge0,\quad \alpha\ge0,\quad \alpha x=0,\\
(x_2^\star,\alpha_2^\star)&=(0,2).
\end{aligned}
\]
\SolCheck $f_1$的无约束极小点$x=2$可行，故约束不改变它；$f_2$的无约束极小点$x=-1$不可行，且$f_2'(x)=2(x+1)>0$对$x\ge0$成立，故边界$x=0$最优。直接比较与KKT解一致。

\SolAnswer 第一问$(x^\star,\alpha^\star)=(2,0)$，约束非活动；第二问$(x^\star,\alpha^\star)=(0,2)$，约束活动。
\end{solution}

\phantomsection\label{struct:V1-C07-CH10}
\phantomsection\label{struct:V1-C07-CH21}
\phantomsection\label{struct:V1-C07-CH22}
\phantomsection\label{struct:V1-C07-CH23}
\phantomsection\label{struct:V1-C07-CH24}
\phantomsection\label{struct:V1-C07-CH25}
\phantomsection\label{struct:V1-C07-CH26}
\phantomsection\label{struct:V1-C07-CH27}
\begin{geometrybox}
\textbf{几何解释。}凸集内任意两点之间的整条线段都留在集合中；凸函数图像位于任意弦的下方。对凸问题，沿最优点与任意候选点连线移动不会先上升后下降，因此局部谷底不能隐藏另一个更低谷底。
\end{geometrybox}

\phantomsection\label{struct:V1-C07-CH19}
% v2.3.1 figure UID: FIG-P109-01
% Proposed post-v2.3.1 polish for FIG-P109-01: protect key-label size and stroke hierarchy.
\tikzset{
  slfig-FIG-P109-01/.style={font=\fontsize{9.2pt}{11.2pt}\selectfont,
    every path/.append style={line cap=round,line join=round}},
  slfig-FIG-P109-01-label/.style={font=\fontsize{9.2pt}{11.2pt}\selectfont,
    fill=white,fill opacity=.92,text opacity=1,inner sep=1.2pt},
  slfig-FIG-P109-01-note/.style={slfig note,rounded corners=1.4mm,
    align=center,font=\fontsize{9.2pt}{11.2pt}\selectfont,
    inner xsep=6pt,inner ysep=4pt,line width=.55pt}
}
\pgfkeysifdefined{/tikz/alt}{}{\tikzset{alt/.code={}}}
\begin{figure}[H]
\centering
\begin{tikzpicture}[slfig-FIG-P109-01,x=1.35cm,y=1.35cm,
  alt={对象—关系—结论：凸集中任意两点的线段仍位于可行域内}]
  \path[fill=SLBlueBG,draw=SLBlue,line width=.82pt]
    plot[smooth cycle,tension=.82] coordinates {
      (-3.7,-.6) (-2.7,-1.45) (-.5,-1.65) (2.0,-1.25)
      (3.55,-.15) (3.2,1.45) (1.2,2.15) (-1.3,2.05) (-3.25,1.15)};
  \coordinate (x) at (-2.35,-.62);
  \coordinate (y) at (2.42,1.02);
  \draw[draw=SLBlue,line width=1.05pt] (x)--(y);
  \fill[SLInk] (x) circle (2.35pt) node[below left=3mm] {$x$};
  \fill[SLInk] (y) circle (2.35pt) node[above right=3mm] {$y$};
  \foreach \f in {.25,.50,.75}{\fill[SLTeal] ($(x)!\f!(y)$) circle (2.25pt);}
  \node[slfig-FIG-P109-01-label,anchor=south east]
    at ($(x)!.50!(y)+(0,3mm)$) {$z=\lambda x+(1-\lambda)y$};
  \node[font=\fontsize{9.2pt}{11.2pt}\selectfont,text=SLInk!78,fill=white,fill opacity=1,text opacity=1,inner sep=1.2pt,anchor=north east]
    at (3.20,1.86) {凸可行域 $C$};
  \node[slfig-FIG-P109-01-note,anchor=north]
    at (0,-1.92) {$x,y\in C,\ \lambda\in[0,1]\ \Longrightarrow\
      \lambda x+(1-\lambda)y\in C$};
\end{tikzpicture}
\caption{凸集中任意两点的线段仍位于可行域内}\label{fig:V1-C07-convex-set}
\end{figure}

用\cref{fig:V1-C07-convex-set}检验候选集合时，不能只找一条留在集合内的线段；定义要求任意$x_1,x_2$和任意$t\in[0,1]$都成立。要否定凸性，只需给出一对端点及一个落到集合外的凸组合。


% KNOWLEDGE-PLAN: KN-V1-C07-BOUNDARY-003 L1
\paragraph{读前自检：凸函数的一阶条件。}先锁定凸函数的一阶条件的条件“可微函数在开凸集上凸”：对$x$完成一项求导、代入或边界比较，所得结果应满足凸函数始终位于其切平面之上。\par
% KNOWLEDGE-PLAN: KN-V1-C07-BOUNDARY-004 L1
\paragraph{读前自检：凸函数的二阶条件。}对凸函数的二阶条件，条件“二次可微函数在开凸集上凸，当且仅当 $\nabla^2f(x)\succeq0$ 对所有$x$成立”不可省；请对$\nabla^2f(x)\succeq0$完成一项求导、代入或边界比较，再用若$f$可微且凸，$C$凸作检查。\par

\SLTeachSection{一阶与二阶凸性条件}{sec:V1-C07-S04}
\knowledgeanchor{PRE-OP-009}{凸函数的一阶条件}
可微函数在开凸集上凸，当且仅当
\[
f(y)\ge f(x)+\nabla f(x)^{\mathsf T}(y-x),\quad \forall x,y.
\]
右侧是$x$处切平面；凸函数始终位于其切平面之上。
\knowledgeanchor{PRE-OP-010}{凸函数的二阶条件}
二次可微函数在开凸集上凸，当且仅当 $\nabla^2f(x)\succeq0$ 对所有$x$成立；若Hessian处处正定，则$f$严格凸。

\phantomsection\label{struct:V1-C07-CH12}
\begin{theorem}[凸问题的一阶最优条件]
若$f$可微且凸，$C$凸，则$x^*\in C$为全局最优点当且仅当
$\nabla f(x^*)^{\mathsf T}(x-x^*)\ge0$ 对所有$x\in C$成立。
\end{theorem}
\phantomsection\label{struct:V1-C07-CH15}
% KNOWLEDGE-PLAN: KN-V1-C07-PROOF-002 L1
\paragraph{读前自检：证明：凸问题的一阶最优条件。}凸问题的一阶最优条件把问题限定在“若条件成立”；请补出$f(x)\ge f(x^*)+\nabla f(x^*)^{\mathsf T}(x-x^*)\ge f(x^*)$到下一关系的一步不等式或求导，检查是否得到单变量函数在$t=0$处右导数不能为负，因此得到所述内积非负。\par

\begin{proof}
若条件成立，凸函数一阶下界给出 $f(x)\ge f(x^*)+\nabla f(x^*)^{\mathsf T}(x-x^*)\ge f(x^*)$。反之若$x^*$最优，对任意$x\in C$，线段$x^*+t(x-x^*)$可行；单变量函数在$t=0$处右导数不能为负，因此得到所述内积非负。
\end{proof}


\SLTeachSection{Jensen不等式}{sec:V1-C07-S05}
\knowledgeanchor{PRE-OP-011}{Jensen不等式}
\cref{pre:V1-C06-finite-jensen}已在信息论首次使用前给出有限加权形式。这里把它放回凸函数主线并证明：若$f$凸，$\lambda_i\ge0$且$\sum_i\lambda_i=1$，则
\[
f\!\left(\sum_i\lambda_ix_i\right)\le\sum_i\lambda_if(x_i).
\]
概率形式为 $f(\mathbb EX)\le\mathbb E f(X)$。若$f$严格凸，等号要求$X$几乎处处为常数；对凹函数不等号方向反转。

\phantomsection\label{struct:V1-C07-CH13}
% KNOWLEDGE-PLAN: KN-V1-C07-LEMMA-001 L1
\paragraph{读前自检：有限Jensen的归纳形式。}在“二点凸性定义可递推推出任意有限个点的加权Jensen不等式”成立时处理有限Jensen的归纳形式：请把“二点凸性定义可递推推出任意有限个点的加权Jensen不等式”中的已声明对象代回有限Jensen的归纳形式的定义，并以对象、条件与结论逐项闭合，且每项都可直接判为成立或失败判定结果。\par

\begin{lemma}[有限Jensen的归纳形式]
二点凸性定义可递推推出任意有限个点的加权Jensen不等式。
\end{lemma}
% KNOWLEDGE-PLAN: KN-V1-C07-PROOF-003 L1
\paragraph{读前自检：证明：有限Jensen的归纳形式。}有限Jensen的归纳形式所用的明确前提是“设最后一个权重为$\lambda_n<1$”；完成补出$z$到下一关系的一步不等式或求导后，核对先对$z,x_n$应用二点凸性，再对$z$的$n-1$点表达应用归纳假设，合并权重后得到结论。\par

\begin{proof}
设最后一个权重为$\lambda_n<1$，把前$n-1$项归一化为点 $z=(1-\lambda_n)^{-1}\sum_{i<n}\lambda_ix_i$。先对$z,x_n$应用二点凸性，再对$z$的$n-1$点表达应用归纳假设，合并权重后得到结论。边界权重为0或1时直接删去零权重点。
\end{proof}


% KNOWLEDGE-PLAN: KN-V1-C07-CONCEPT-008 L1
\paragraph{读前自检：拉格朗日乘子法。}等式约束$h(x)=0$在候选最优点$x^*$处满足约束梯度正则；请核对$\nabla f(x^*)+J_h(x^*)^{\mathsf T}\beta=0$，并说明乘子$\beta$给出约束右端微扰下最优值的一阶敏感度。\par

\SLAdvancedSection{拉格朗日乘子与拉格朗日函数}{sec:V1-C07-S06}
\knowledgeanchor{PRE-OP-012}{拉格朗日乘子法}
等式约束下，若约束梯度正则，最优点处目标梯度属于约束梯度张成空间。乘子刻画放宽约束对最优值的一阶敏感度。
\knowledgeanchor{PRE-OP-013}{拉格朗日函数}
对标准形式定义
\[
L(x,\alpha,\beta)=f(x)+\sum_{i=1}^{m}\alpha_i c_i(x)+\sum_{j=1}^{r}\beta_jh_j(x),
\qquad \alpha_i\ge0.
\]
不等式乘子非负保证对可行点有 $L(x,\alpha,\beta)\le f(x)$；等式项在可行点为零。


% KNOWLEDGE-PLAN: KN-V1-C07-CONCEPT-009 L1
\paragraph{读前自检：原问题。}原问题以$f(x)$及$c_i(x)\le0,h_j(x)=0$定义可行域，最优值记为$p^*$；请先判断一个候选$x$是否可行，再核对可行域为空时按约定$p^*=+\infty$。\par
% KNOWLEDGE-PLAN: KN-V1-C07-CONCEPT-010 L1
\paragraph{读前自检：拉格朗日对偶性。}拉格朗日对偶性从$L(x,\alpha,\beta)=f(x)+\sum_i\alpha_ic_i(x)+\sum_j\beta_jh_j(x)$出发；请在$\alpha\ge0$下对$x$取下确界，核对$g(\alpha,\beta)$不超过任一原可行点的目标值。\par
% KNOWLEDGE-PLAN: KN-V1-C07-CONCEPT-011 L1
\paragraph{读前自检：原始问题。}原始问题采用最小化$f(x)$并满足$c_i(x)\le0,h_j(x)=0$的标准形；请把$a(x)\ge b$改写成$b-a(x)\le0$，核对可行集合保持不变。\par
% KNOWLEDGE-PLAN: KN-V1-C07-PREREQUISITE-010 L1
\paragraph{读前自检：对偶函数。}从凸优化与拉格朗日对偶中的对偶函数的条件“$g(\alpha,\beta)=\inf_x L(x,\alpha,\beta).$”出发，请把$g(\alpha,\beta)$展开一次并检查其类型或边界；所得量应符合无论原问题是否凸，$g$总是凹函数。\par
% KNOWLEDGE-PLAN: KN-V1-C07-CONCEPT-012 L1
\paragraph{读前自检：对偶问题。}对偶问题在$\alpha\ge0$、$\beta$自由的域上最大化$g(\alpha,\beta)$；请对任一对偶可行乘子写出$g(\alpha,\beta)\le p^*$，再核对最优值$d^*$仍是原最优值的下界。\par
% KNOWLEDGE-PLAN: KN-V1-C07-CONCEPT-013 L1
\paragraph{读前自检：对偶问题。}对偶问题$\max_{\alpha\ge0,\beta}g(\alpha,\beta)$区分不等式与等式乘子；请逐项核对$\alpha$的非负域、$\beta$的自由域，并确认最大值记作$d^*$。\par

\SLAdvancedSection{原问题、对偶函数与对偶问题}{sec:V1-C07-S07}
\knowledgeanchor{PRE-OP-014}{原问题}
\knowledgeanchor{SLM-C-00-001}{拉格朗日对偶性}
\knowledgeanchor{SLM-C-01-001}{原始问题}
原问题就是本章标准约束最小化，其最优值记为$p^*$。可用广义拉格朗日函数把不可行点惩罚为$+\infty$，从而与原问题等价。
\knowledgeanchor{PRE-OP-015}{对偶函数}
\[
g(\alpha,\beta)=\inf_x L(x,\alpha,\beta).
\]
它是关于乘子的仿射函数族的逐点下确界，因此无论原问题是否凸，$g$总是凹函数。
\knowledgeanchor{PRE-OP-016}{对偶问题}
\knowledgeanchor{SLM-C-02-001}{对偶问题}
对偶问题为 $\max_{\alpha\ge0,\beta}g(\alpha,\beta)$，最优值记为$d^*$。对偶可行乘子给出原最优值的下界。


% KNOWLEDGE-PLAN: KN-V1-C07-CONCEPT-014 L1
\paragraph{读前自检：弱对偶。}先锁定凸优化与拉格朗日对偶中的弱对偶的条件“对任何原可行$x$与对偶可行$(\alpha,\beta)$”：代入正文给出的取等边界检验$g(\alpha,\beta)\le L(x,\alpha,\beta)\le f(x),$，所得结果应满足不等号方向、边界取等和定义域彼此相容。\par
% KNOWLEDGE-PLAN: KN-V1-C07-CONCEPT-015 L1
\paragraph{读前自检：原始问题和对偶问题的关系。}对原始问题和对偶问题的关系，条件“对任何原可行$x$与对偶可行$(\alpha,\beta)$”不可省；请计算$g(\alpha,\beta)\le L(x,\alpha,\beta)\le f(x),$两边之差并判号，再用不等号方向、边界取等和定义域彼此相容作检查。\par
% KNOWLEDGE-PLAN: KN-V1-C07-THEOREM-003 L1
\paragraph{读前自检：定理C 1。}定理C 1把问题限定在“对任何原可行$x$与对偶可行$(\alpha,\beta)$”；请把$g(\alpha,\beta)\le L(x,\alpha,\beta)\le f(x),$用到的关键定义展开并完成下一步变形，检查是否得到变形没有增加前提且得到证明所述结论。\par

\SLAdvancedSection{弱对偶、强对偶与Slater条件}{sec:V1-C07-S08}
\knowledgeanchor{PRE-OP-017}{弱对偶}
\knowledgeanchor{SLM-C-03-001}{原始问题和对偶问题的关系}
\knowledgeanchor{SLM-C-01-002}{定理C.1}
对任何原可行$x$与对偶可行$(\alpha,\beta)$，有
\[
g(\alpha,\beta)\le L(x,\alpha,\beta)\le f(x),
\]
故 $d^*\le p^*$。差$p^*-d^*$称为对偶间隙。

% KNOWLEDGE-PLAN: KN-V1-C07-THEOREM-004 L1
\paragraph{读前自检：推论C 1。}在“若找到原可行与对偶可行解且目标值相等”成立时处理推论C 1：请把“若找到原可行与对偶可行解且目标值相等”中的已声明对象代回推论C 1的定义，并以两者分别为各自问题的最优解判定结果。\par
% KNOWLEDGE-PLAN: KN-V1-C07-CONCEPT-016 L1
\paragraph{读前自检：强对偶。}凸优化与拉格朗日对偶中的强对偶所用的明确前提是“强对偶指 $d^*=p^*$”；完成检查$d^*=p^*$中每项的类型后完成等式变形后，核对它不是所有约束问题都自动具有的性质。\par

\knowledgeanchor{SLM-C-01-003}{推论C.1}
若找到原可行与对偶可行解且目标值相等，则两者分别为各自问题的最优解。
\knowledgeanchor{PRE-OP-018}{强对偶}
强对偶指 $d^*=p^*$；它不是所有约束问题都自动具有的性质。
\knowledgeanchor{PRE-OP-019}{Slater条件}
考虑有限维凸问题：目标与不等式函数凸，等式约束写成$Ax=b$，且各函数有共同定义域。若存在$\widetilde x\in\operatorname{ri}(\operatorname{dom}f_0\cap\cdots\cap\operatorname{dom}f_m)$满足$f_i(\widetilde x)<0$（所有非仿射不等式）并满足$A\widetilde x=b$，则称满足Slater条件；它保证强对偶，并在常见可达性条件下保证对偶最优乘子存在。

% KNOWLEDGE-PLAN: KN-V1-C07-THEOREM-005 L1
\paragraph{读前自检：定理C 2。}定理C.2的强对偶证明用超平面分离原问题可达集合与更优不可达点；请列出集合非空、凸性及相对内部或闭包边界条件，并指出缺少这些条件时分离步骤不能直接调用。\par

\knowledgeanchor{SLM-C-02-002}{定理C.2}
\begin{theorem}[凸问题的强对偶]
有限维凸目标与凸不等式约束、仿射等式约束在上述相对内部Slater条件下满足 $d^*=p^*$。
\end{theorem}

\begin{keypointbox}[title=扩展讨论：分离定理在这里做什么]
强对偶证明常用有限维凸集分离：主线只需把它理解为“把原问题可达集合与更优但不可达的点用超平面隔开”。严格版本必须说明集合非空、凸性以及相对内部或闭包边界条件；这些条件不作为本章计算KKT残差的前置要求。
\end{keypointbox}

% DERIVATION-CHECK: step-1; step-2; step-3
\phantomsection\label{struct:V1-C07-CH11}
% KNOWLEDGE-PLAN: KN-V1-C07-ALGORITHM_IDEA-001 L2

% KNOWLEDGE-PLAN: KN-V1-C07-ALGORITHM_IDEA-002 L2
\begin{keypointbox}[title={候选原--对偶解的KKT残差检验：五步阅读}]
\textbf{输入。}写出数据对象、固定参数与必须成立的条件。\par
\textbf{初始化。}标明主状态、计数器和第一个合法状态。\par
\textbf{核心更新。}只手算一次从旧状态到候选状态的更新，并指出何时提交。\par
\textbf{数学停止。}写出能直接核验的停止证书；预算耗尽不等同于收敛。\par
\textbf{输出。}列出返回对象、实际计数、中心状态和用于复核的诊断。
\end{keypointbox}

\begin{AlgorithmContract}{
  input={候选原变量$x$、乘子$\alpha,\beta$、约束/目标函数及梯度、容差$\varepsilon$与范数规则},
  output={最后合法残差记录、活动集、通过标志、实际约束检查数$c_{\rm done}$、状态与最终诊断},
  preconditions={候选与乘子维数匹配，$\varepsilon>0$，所有声明函数接口可在候选点调用},
  initialization={置残差记录和活动集为空、$R=0,c_{\rm done}=0$，先形成驻点残差候选},
  loop={按固定下标扫描$m$个不等式与$r$个等式约束，逐项形成可行性、乘子与互补残差},
  update={单项函数值和残差有限后才原子追加记录、更新活动集与最大残差},
  domain={所有函数值、梯度、乘子和残差为有限实数；活动判定使用预先给定容差},
  failure={维数、容差或接口非法返回$\mathtt{invalid\_input}$；函数/梯度或残差非有限返回$\mathtt{numerical\_failure}$并保留最后合法残差前缀},
  stop={驻点项和全部$m+r$个约束核验完成，或首次失败时停止},
  budget={恰核验一个驻点残差、$m$个不等式与$r$个等式，约束检查上限$m+r$},
  status={$\mathtt{completed}$、$\mathtt{invalid\_input}$、$\mathtt{numerical\_failure}$},
  iterations={$c_{\rm done}$计实际完成的约束项，有限核验不使用优化迭代数冒充收敛},
  diagnostics={五类最大残差、主导约束、活动集、通过标志、失败约束下标与求值阶段},
  complexity={稠密线性约束核验为$O((m+r)d)$时间，残差和活动集额外空间$O(m+r)$}
}
\begin{algorithm}[H]
\caption{候选原--对偶解的KKT残差检验}\label{alg:V1-C07-kkt-check}
\KwIn{候选$x,\alpha,\beta$；函数与梯度；容差$\varepsilon>0$}
\KwOut{最后合法残差、活动集、通过标志；$c_{\rm done}$；$\mathtt{status}$与诊断}
若维数、容差、范数或接口非法，返回空证书、$c_{\rm done}=0$、$\mathtt{invalid\_input}$及字段诊断\;
初始化$\mathcal A\leftarrow\varnothing,R\leftarrow0,c_{\rm done}\leftarrow0$，形成$R_{\mathrm s}^+\leftarrow\|\nabla_xL(x,\alpha,\beta)\|$\;
若驻点量非有限，返回空证书、$\mathtt{numerical\_failure}$及梯度诊断；否则提交$R_{\mathrm s}\leftarrow R_{\mathrm s}^+,R\leftarrow R_{\mathrm s}$\;
\For{$i\leftarrow1$ \KwTo $m$}{
  形成第$i$项原可行、对偶可行与互补残差候选；非有限时返回旧记录、$\mathtt{numerical\_failure}$及$i$诊断\;
  原子追加第$i$项残差、更新$R$与活动集并增加$c_{\rm done}$\;
}
\For{$j\leftarrow1$ \KwTo $r$}{形成$|h_j(x)|$；非有限时返回旧记录、$\mathtt{numerical\_failure}$及$j$诊断；否则原子追加、更新$R$并增加$c_{\rm done}$\;}
令通过标志为$\mathbf1\{R\le\varepsilon\}$；返回完整证书、$\mathtt{completed}$及主导残差/容差最终诊断\;
\end{algorithm}
\end{AlgorithmContract}
\SLRobustImplementationStrip{首次阅读只计算驻点、原可行、对偶可行、互补残差并取最大值；实现时再检查非有限值、活动集和返回状态}

\textbf{收敛与非收敛说明。}KKT残差检验只评价一个给定候选点，不自行生成参数序列，因此不讨论检验器自身的迭代收敛。返回“未通过”表示候选点没有满足容差，不等同于某个外部优化器已经被证明不收敛。


\phantomsection\label{struct:V1-C07-CH17}
\begin{probabilitybox}
归一化约束、概率非负约束和对数似然目标都可纳入KKT框架。非负乘子可解释为某个概率边界对最优目标的影子价格，但只有在模型和约束资格成立时该解释才可靠。
\end{probabilitybox}
\phantomsection\label{struct:V1-C07-CH18}
\begin{optimizationbox}
\textbf{优化解释。}对偶变量为原问题提供下界和敏感度；原--对偶间隙给出可计算停止准则。KKT把可行性、驻点和活动约束统一成方程与不等式系统，是后续支持向量机、最大熵和约束统计估计的共同语言。
\end{optimizationbox}

\phantomsection\label{struct:V1-C07-CH28}
\phantomsection\label{struct:V1-C07-CH29}
% KNOWLEDGE-PLAN: KN-V1-C07-BOUNDARY-006 L1
\paragraph{读前自检：复杂度分析：弱对偶、强对偶与Slater条件。}弱对偶、强对偶与Slater条件中的“采用稀疏或矩阵自由方法时，复杂度应以非零元和迭代次数参数化”决定可执行范围；请由$(m+r)d$的总成本剥离一次迭代或一次样本因子，结果须满足稀疏结构、Schur补和迭代法可降低实际成本。\par

\begin{complexitybox}
\textbf{复杂度假设。}以下界假设函数值和梯度已经可用，稠密线性约束矩阵按全部$(m+r)d$个系数访问，稠密KKT系统用直接分解求解；函数求值、自动微分以及稀疏预处理成本需另计。采用稀疏或矩阵自由方法时，复杂度应以非零元和迭代次数参数化。
设$d$为变量维数，$m,r$分别为不等式与等式个数。\textbf{核验复杂度：}KKT残差检验只核验给定候选点，不更新模型参数。一次稠密线性约束核验为$O((m+r)d)$，直接求解稠密KKT线性系统则通常为$O((d+r)^3)$。\textbf{预测复杂度：}检验器不定义新样本预测映射。\textbf{存储复杂度：}保存稠密约束矩阵为$O((m+r)d)$，残差与活动集额外为$O(m+r)$；稀疏结构、Schur补和迭代法可降低实际成本。
\end{complexitybox}

\phantomsection\label{struct:V1-C07-CH30}
\phantomsection\label{struct:V1-C07-CH31}
\begin{applicabilitybox}
\textbf{适用条件：}凸优化理论适用于凸目标、凸不等式和仿射等式；Slater条件常用来保证强对偶，KKT可作为充分必要证书。\textbf{局限性：}非凸问题中KKT通常只给必要候选，约束资格失败时最优点甚至可能没有乘子；严格凸只保证解至多唯一，不自动保证存在；数值算法还受条件数与尺度影响。
\end{applicabilitybox}

\phantomsection\label{struct:V1-C07-CH32}
% KNOWLEDGE-PLAN: KN-V1-C07-BOUNDARY-008 L1
\paragraph{读前自检：常见错误与误用边界：弱对偶、强对偶与Slater条件。}把弱对偶、强对偶与Slater条件放在“边界框首项禁止把局部最优自动当全局最优而未检查凸性”下考察：把固定错误“把局部最优自动当全局最优而未检查凸性”中的对象、操作与结论按本节定义拆开，指出第一个不满足的定义条件，再把该处改为本节允许的写法，并检查改写后的运算或判断有定义，且不再触发“把局部最优自动当全局最优而未检查凸性”。\par

\begin{warningbox}
典型误用包括：混淆最大化与最小化符号；把 $c_i\ge0$ 直接配非负乘子却仍套用本章公式；认为Hessian在一个点半正定即可证明全域凸；把局部最优自动当全局最优而未检查凸性；把弱对偶误写成$d^*\ge p^*$；漏掉KKT中的原可行、对偶可行或互补松弛。
\end{warningbox}

\phantomsection\label{struct:V1-C07-CH33}
\begin{comparisonbox}
\textbf{概念对比：}凸集是域的几何性质，凸函数是函数沿线段的性质；严格凸关乎解唯一性，强凸还给出统一曲率下界。弱对偶总成立，强对偶需要附加条件。活动约束满足等号，正乘子必对应活动约束，但活动约束可有零乘子。KKT是最优性系统，拉格朗日函数本身不是无条件的惩罚目标。
\end{comparisonbox}

\begin{example}[原点到半空间的KKT投影]\label{exm:V1-C07-halfspace-kkt}
设$a\in\mathbb R^d\setminus\{0\}$、$b>0$。求$\min_x\tfrac12\|x\|_2^2$且$a^{\mathsf T}x\ge b$的KKT解。
\end{example}
\SLExampleSolutionHeading{exm:V1-C07-halfspace-kkt}
\begin{solution}
\SLReadTranslation{求原点到半空间$a^{\mathsf T}x\ge b$的最小范数点，同时给出KKT乘子和最优目标值。}
\SolGiven $a\in\mathbb R^d\setminus\{0\}$、$b>0$；目标严格凸，约束仿射且存在严格可行点，因此KKT解存在并唯一确定原变量。
\SLMethodTrigger{把约束写成$b-a^{\mathsf T}x\le0$，由驻点条件与互补松弛解出候选；再用Cauchy--Schwarz给出独立范数下界。}
\SolPlan 写出拉格朗日函数与KKT条件，由$x=\alpha a$判定约束活动并求$\alpha,x$；最后验证候选达到所有可行点的范数下界。
\SolDerive
目标严格凸，标准约束$c(x)=b-a^{\mathsf T}x\le0$仿射，可行域为非空闭半空间；取足够大的$x=ta$可严格可行，故Slater成立。拉格朗日函数与KKT条件为
\[
\begin{aligned}
L(x,\alpha)&=\frac12\|x\|_2^2+\alpha(b-a^{\mathsf T}x),\\
x-\alpha a&=0,\qquad
b-a^{\mathsf T}x\le0,\qquad
\alpha\ge0,\\
\alpha(b-a^{\mathsf T}x)&=0.
\end{aligned}
\]
由$x=\alpha a$且$x=0$不可行，约束必须活动。因此
\[
\begin{aligned}
a^{\mathsf T}x^\star=b
&\Longrightarrow \alpha^\star\|a\|_2^2=b,\\
\alpha^\star&=\frac{b}{\|a\|_2^2}>0,\qquad
x^\star=\frac{b\,a}{\|a\|_2^2}.
\end{aligned}
\]
目标严格凸而约束集凸，故该KKT点是唯一全局最优解。
\SolCheck 任意可行$x$由Cauchy--Schwarz不等式满足$b\le a^{\mathsf T}x\le\|a\|_2\|x\|_2$，所以$\|x\|_2\ge b/\|a\|_2$；候选$x^\star=ba/\|a\|_2^2$恰好取等。

\SolAnswer $x^\star=ba/\|a\|_2^2$，$\alpha^\star=b/\|a\|_2^2$，最优值为$b^2/(2\|a\|_2^2)$；该点是原点到半空间边界的正交投影。
\end{solution}

\begin{chapterexercisebox}
\phantomsection\label{struct:V1-C07-CH34}
\phantomsection\label{struct:V1-C07-CH35}
本章共18道练习。每道题干后立即给出同号解析；长解析可自然跨页，续页标题保留题号。
\end{chapterexercisebox}

\begin{SLChapterExercisePairSection}

\Needspace{6\baselineskip}
\begin{exercise}\label{ex:V1-C07-S01-01}
把“最大化 $a^{\mathsf T}x$，满足 $Bx\ge b$ 与 $\mathbf1^{\mathsf T}x=1$”改写成本章标准形式。
\end{exercise}
\phantomsection\label{sol:V1-C07-S01-01}
\begin{chapterexercisesolution}{ex:V1-C07-S01-01}
设$x\in\mathbb R^d$、$a\in\mathbb R^d$、$B\in\mathbb R^{m\times d}$、$b\in\mathbb R^m$。原可行域与标准最小化形式为
\[
\begin{aligned}
\mathcal C
&=\{x\in\mathbb R^d:Bx\ge b,\ \mathbf1^{\mathsf T}x=1\},\\
\min_{x\in\mathbb R^d}\quad
f(x)&=-a^{\mathsf T}x,\\
\text{s.t.}\quad
c_i(x)&=b_i-B_i x\le0,\quad i=1,\ldots,m,\\
h(x)&=\mathbf1^{\mathsf T}x-1=0.
\end{aligned}
\]
若$\mathcal C\ne\varnothing$，该改写保持可行点和最优解。各目标与约束均为仿射函数，所以任一满足$\alpha_i\ge0$的KKT点都足以证明全局最优；若要保证最优点必存在相应乘子，还需适当约束资格条件。
\end{chapterexercisesolution}

\Needspace{6\baselineskip}
\begin{exercise}\label{ex:V1-C07-S01-02}
对岭回归 $\min_w\|Xw-y\|_2^2+\lambda\|w\|_2^2$，指出变量、目标和参数，并说明是否有显式约束。
\end{exercise}
\phantomsection\label{sol:V1-C07-S01-02}
\begin{chapterexercisesolution}{ex:V1-C07-S01-02}
声明$X\in\mathbb R^{n\times d}$、$y\in\mathbb R^n$、变量$w\in\mathbb R^d$与给定参数$\lambda\ge0$。目标和可行域为
\[
\begin{aligned}
f(w)&=\|Xw-y\|_2^2+\lambda\|w\|_2^2,\\
\mathcal C&=\mathbb R^d\qquad\text{（无显式约束）},\\
\nabla f(w)&=2X^{\mathsf T}(Xw-y)+2\lambda w,\\
\nabla^2f(w)&=2(X^{\mathsf T}X+\lambda I_d)\succeq0.
\end{aligned}
\]
无约束KKT退化为$\nabla f(w^\star)=0$。若$\lambda>0$，Hessian正定，故唯一全局解存在；若$\lambda=0$，还需$X$满列秩才保证唯一。
\end{chapterexercisesolution}

\Needspace{6\baselineskip}
\begin{exercise}\label{ex:V1-C07-S02-01}
求 $f(x)=x^3-3x$ 在区间$[-2,2]$上的全部局部与全局最优点。
\end{exercise}
\phantomsection\label{sol:V1-C07-S02-01}
\begin{chapterexercisesolution}{ex:V1-C07-S02-01}
目标$f(x)=x^3-3x$连续，可行域$\mathcal C=[-2,2]$非空紧，故全局极值存在。内部驻点与单调性为
\[
\begin{aligned}
f'(x)&=3(x^2-1)=0
\Longleftrightarrow x\in\{-1,1\},\\
f'(x)&>0\ \text{于 }(-2,-1)\cup(1,2),\qquad
f'(x)<0\ \text{于 }(-1,1).
\end{aligned}
\]
因此相对于可行域的局部分类及函数值为
\[
\begin{array}{c|cccc}
x&-2&-1&1&2\\ \hline
\text{局部类型}&\text{极小}&\text{极大}&\text{极小}&\text{极大}\\
f(x)&-2&2&-2&2
\end{array}
\]
故全局极小点为$\{-2,1\}$，全局极大点为$\{-1,2\}$。端点采用单侧可行邻域，也属于局部最优点。
\end{chapterexercisesolution}

\Needspace{6\baselineskip}
\begin{exercise}\label{ex:V1-C07-S02-02}
证明若可行域非空且紧、目标连续，则全局最优解存在。
\end{exercise}
\phantomsection\label{sol:V1-C07-S02-02}
\begin{chapterexercisesolution}{ex:V1-C07-S02-02}
设可行域$\mathcal C\subseteq\mathbb R^d$非空且紧，目标$f:\mathcal C\to\mathbb R$连续。连续映射保持紧性，因此
\[
\begin{aligned}
\mathcal C\ne\varnothing,\ \mathcal C\text{紧}
&\Longrightarrow f(\mathcal C)\ne\varnothing,\ f(\mathcal C)\text{紧},\\
m&:=\inf f(\mathcal C)\in f(\mathcal C),\\
\exists x^\star\in\mathcal C:\quad
f(x^\star)&=m=\min_{x\in\mathcal C}f(x).
\end{aligned}
\]
这就是Weierstrass定理；该存在性结论不需要凸性或KKT。若可行域不闭或无界，下确界可能不被任何可行点取得。
\end{chapterexercisesolution}

\Needspace{6\baselineskip}
\begin{exercise}\label{ex:V1-C07-S03-01}
证明半空间 $C=\{x:a^{\mathsf T}x\le b\}$ 为凸集。
\end{exercise}
\phantomsection\label{sol:V1-C07-S03-01}
\begin{chapterexercisesolution}{ex:V1-C07-S03-01}
对象为半空间$C=\{x\in\mathbb R^d:a^{\mathsf T}x\le b\}$。任取$x,y\in C$及$t\in[0,1]$，
\[
\begin{aligned}
a^{\mathsf T}(tx+(1-t)y)
&=t\,a^{\mathsf T}x+(1-t)a^{\mathsf T}y\\
&\le tb+(1-t)b=b,\\
tx+(1-t)y&\in C.
\end{aligned}
\]
故$C$为凸集。证明也覆盖$a=0$时的退化空集或全空间；本题只判断集合凸性，不涉及目标函数或KKT。
\end{chapterexercisesolution}

\Needspace{6\baselineskip}
\begin{exercise}\label{ex:V1-C07-S03-02}
判断集合 $\{x\in\mathbb R^2:\|x\|_2=1\}$ 是否凸，并给出反例点。
\end{exercise}
\phantomsection\label{sol:V1-C07-S03-02}
\begin{chapterexercisesolution}{ex:V1-C07-S03-02}
集合$C=\{x\in\mathbb R^2:\|x\|_2=1\}$只有圆周而不含内部。取两个可行点及$t=1/2$，
\[
\begin{aligned}
x&=(1,0)^{\mathsf T}\in C,\qquad
y=(-1,0)^{\mathsf T}\in C,\\
\frac12x+\frac12y&=(0,0)^{\mathsf T},\\
\left\|\frac12x+\frac12y\right\|_2&=0\ne1.
\end{aligned}
\]
故$C$不凸。相反，$\|x\|_2\le1$定义的闭球因三角不等式而凸；集合判定本身不适用KKT。
\end{chapterexercisesolution}

\Needspace{6\baselineskip}
\begin{exercise}\label{ex:V1-C07-S04-01}
用Hessian判断 $f(x)=\tfrac12x^{\mathsf T}Qx+b^{\mathsf T}x$ 的凸性与严格凸性。
\end{exercise}
\phantomsection\label{sol:V1-C07-S04-01}
\begin{chapterexercisesolution}{ex:V1-C07-S04-01}
对象为$f:\mathbb R^d\to\mathbb R$，$Q\in\mathbb R^{d\times d}$、$b\in\mathbb R^d$，定义域是凸开集$\mathbb R^d$。令$S=(Q+Q^{\mathsf T})/2$，则
\[
\begin{aligned}
x^{\mathsf T}Qx&=x^{\mathsf T}Sx,\\
\nabla f(x)&=Sx+b,\\
\nabla^2f(x)&=S,\\
f\text{凸}&\Longleftrightarrow S\succeq0,\qquad
f\text{严格凸}\Longleftrightarrow S\succ0.
\end{aligned}
\]
若$S$有负特征值，沿对应方向二阶曲率为负，故不凸。无约束最优化时KKT仅为$Sx+b=0$；解的存在唯一性还取决于$S$与$b$。
\end{chapterexercisesolution}

\Needspace{6\baselineskip}
\begin{exercise}\label{ex:V1-C07-S04-02}
证明 $f(x)=\log\sum_{k=1}^{K}e^{x_k}$ 是凸函数。
\end{exercise}
\phantomsection\label{sol:V1-C07-S04-02}
\begin{chapterexercisesolution}{ex:V1-C07-S04-02}
函数$f:\mathbb R^K\to\mathbb R$处处有限。令$Z=\sum_{j=1}^Ke^{x_j}>0$及$p_k=e^{x_k}/Z$，则$\boldsymbol p$是有限支持上的概率向量。逐次求导，
\[
\begin{aligned}
\nabla f(x)&=\boldsymbol p,\\
\nabla^2f(x)&=\operatorname{diag}(\boldsymbol p)
-\boldsymbol p\boldsymbol p^{\mathsf T},\\
v^{\mathsf T}\nabla^2f(x)v
&=\sum_{k=1}^Kp_kv_k^2
-\left(\sum_{k=1}^Kp_kv_k\right)^2\\
&=\operatorname{Var}_{K\sim\boldsymbol p}(v_K)\ge0.
\end{aligned}
\]
故$f$凸；沿$\mathbf1$方向有$f(x+t\mathbf1)=f(x)+t$，所以在整个$\mathbb R^K$上不严格凸。本题是函数凸性证明，无独立约束，KKT不另行使用。
\end{chapterexercisesolution}

\Needspace{6\baselineskip}
\begin{exercise}\label{ex:V1-C07-S05-01}
用Jensen证明正数 $x_1,\ldots,x_n$ 的算术平均不小于几何平均。
\end{exercise}
\phantomsection\label{sol:V1-C07-S05-01}
\begin{chapterexercisesolution}{ex:V1-C07-S05-01}
给定$n\ge1$与$x_i>0$，算术平均和对数均有定义。凹函数$\log$的Jensen不等式给出
\[
\begin{aligned}
\log\left(\frac1n\sum_{i=1}^nx_i\right)
&\ge\frac1n\sum_{i=1}^n\log x_i\\
&=\log\left(\prod_{i=1}^nx_i\right)^{1/n},\\
\frac1n\sum_{i=1}^nx_i
&\ge\left(\prod_{i=1}^nx_i\right)^{1/n}.
\end{aligned}
\]
指数函数严格递增，故保持不等号；因权重均为正，等号当且仅当$x_1=\cdots=x_n$。
\end{chapterexercisesolution}

\Needspace{6\baselineskip}
\begin{exercise}\label{ex:V1-C07-S05-02}
设$X>0$且相关期望有限，用Jensen比较 $E[1/X]$ 与 $1/E[X]$。
\end{exercise}
\phantomsection\label{sol:V1-C07-S05-02}
\begin{chapterexercisesolution}{ex:V1-C07-S05-02}
假定$X>0$几乎处处、$0<\mathbb E[X]<\infty$且$\mathbb E[1/X]<\infty$，使两侧均为有限实数。函数$\phi(x)=1/x$在$(0,\infty)$上满足
\[
\begin{aligned}
\phi''(x)&=\frac2{x^3}>0,\\
\phi(\mathbb E[X])
&\le\mathbb E[\phi(X)],\\
\frac1{\mathbb E[X]}&\le\mathbb E\!\left[\frac1X\right].
\end{aligned}
\]
严格凸性的等号条件为$X$几乎处处为常数；若允许$\mathbb E[1/X]=+\infty$，不等式仍可按扩展实数成立，但不再是有限量比较。
\end{chapterexercisesolution}

\Needspace{6\baselineskip}
\begin{exercise}\label{ex:V1-C07-S06-01}
用乘子法求 $\min x_1^2+x_2^2$，满足 $x_1+x_2=1$。
\end{exercise}
\phantomsection\label{sol:V1-C07-S06-01}
\begin{chapterexercisesolution}{ex:V1-C07-S06-01}
目标$f(x)=x_1^2+x_2^2$严格凸，可行域$\mathcal C=\{x\in\mathbb R^2:x_1+x_2=1\}$非空且由仿射等式定义。拉格朗日函数及等式KKT条件为
\[
\begin{aligned}
L(x,\beta)&=x_1^2+x_2^2+\beta(x_1+x_2-1),\\
\nabla_xL&=\begin{bmatrix}2x_1+\beta\\2x_2+\beta\end{bmatrix}
=\boldsymbol0,\qquad x_1+x_2=1,\\
x_1=x_2&=-\frac\beta2,\qquad
x_1^\star=x_2^\star=\frac12,\quad \beta^\star=-1,\\
f(x^\star)&=\frac14+\frac14=\frac12.
\end{aligned}
\]
仿射等式与严格凸目标使该KKT点成为唯一全局最优解。
\end{chapterexercisesolution}

\Needspace{6\baselineskip}
\begin{exercise}\label{ex:V1-C07-S06-02}
解释标准约束$c(x)\le0$的乘子为何要求非负，并说明若允许负值会破坏什么下界关系。
\end{exercise}
\phantomsection\label{sol:V1-C07-S06-02}
\begin{chapterexercisesolution}{ex:V1-C07-S06-02}
考虑原问题$\min_{x\in\mathcal C}f(x)$，标准不等式$c_i(x)\le0$，拉格朗日函数
$L(x,\alpha)=f(x)+\sum_i\alpha_ic_i(x)$。对任意原可行$x$，若$\alpha_i\ge0$，则
\[
\begin{aligned}
\alpha_ic_i(x)&\le0,\\
L(x,\alpha)&\le f(x),\\
g(\alpha):=\inf_zL(z,\alpha)
&\le L(x,\alpha)\le f(x),\\
g(\alpha)&\le p^\star.
\end{aligned}
\]
这就是弱对偶下界。若允许某个$\alpha_i<0$，则在$c_i(x)<0$处可能有$\alpha_ic_i(x)>0$，中间的$L(x,\alpha)\le f(x)$失效，故KKT的对偶可行性必须要求$\alpha\ge0$。
\end{chapterexercisesolution}

\Needspace{6\baselineskip}
\begin{exercise}\label{ex:V1-C07-S07-01}
对无约束二次函数 $L(x,\alpha)=\tfrac12x^2+\alpha(1-x)$，求对偶函数。
\end{exercise}
\phantomsection\label{sol:V1-C07-S07-01}
\begin{chapterexercisesolution}{ex:V1-C07-S07-01}
对偶函数的定义域先取所有$\alpha\in\mathbb R$，并对无约束变量$x\in\mathbb R$取下确界。因$x$方向二阶导数为1，极小点存在且唯一：
\[
\begin{aligned}
g(\alpha)
&=\inf_{x\in\mathbb R}\left\{\frac12x^2+\alpha(1-x)\right\},\\
\frac{\partial L}{\partial x}&=x-\alpha=0
\Longrightarrow x_\alpha=\alpha,\\
g(\alpha)&=L(\alpha,\alpha)
=\frac12\alpha^2+\alpha(1-\alpha)
=\alpha-\frac12\alpha^2,\\
g''(\alpha)&=-1<0.
\end{aligned}
\]
若$\alpha$对应标准不等式$1-x\le0$，对偶可行域还须限制为$\alpha\ge0$。
\end{chapterexercisesolution}

\Needspace{6\baselineskip}
\begin{exercise}\label{ex:V1-C07-S07-02}
证明对偶函数作为仿射函数的逐点下确界是凹函数。
\end{exercise}
\phantomsection\label{sol:V1-C07-S07-02}
\begin{chapterexercisesolution}{ex:V1-C07-S07-02}
设乘子域$\mathcal D$为凸集，且对每个固定$x$，$L(x,u)$关于$u$仿射。先取$g(u),g(v)$有限的$u,v\in\mathcal D$及$t\in[0,1]$，
\[
\begin{aligned}
L(x,tu+(1-t)v)
&=tL(x,u)+(1-t)L(x,v)\\
&\ge tg(u)+(1-t)g(v),\qquad \forall x,\\
g(tu+(1-t)v)
&=\inf_xL(x,tu+(1-t)v)\\
&\ge tg(u)+(1-t)g(v).
\end{aligned}
\]
故$g$在$\mathcal D$上凹；若某端值为$-\infty$，扩展实数意义下的不等式平凡成立。该结论不要求原问题凸，但把$g$用作有效下界仍需乘子满足对偶可行性。
\end{chapterexercisesolution}

\Needspace{6\baselineskip}
\begin{exercise}\label{ex:V1-C07-S08-01}
若某对偶可行点给出下界7，某原可行点目标为9，可以对最优值和对偶间隙作何判断？
\end{exercise}
\phantomsection\label{sol:V1-C07-S08-01}
\begin{chapterexercisesolution}{ex:V1-C07-S08-01}
设问题为最小化，已知一个对偶可行点的对偶值$g(\alpha)=7$及一个原可行点的目标值$f(x)=9$。弱对偶给出
\[
\begin{aligned}
7=g(\alpha)&\le d^\star\le p^\star\le f(x)=9,\\
p^\star&\in[7,9],\\
f(x)-g(\alpha)&=9-7=2,\\
0\le p^\star-d^\star&\le2.
\end{aligned}
\]
没有强对偶或KKT证书时，不能断言任一给定可行点最优；只有上下界相等才直接证明最优。
\end{chapterexercisesolution}

\Needspace{6\baselineskip}
\begin{exercise}\label{ex:V1-C07-S08-02}
检查问题 $\min_x x^2$，满足 $x\le1$ 是否满足Slater条件。
\end{exercise}
\phantomsection\label{sol:V1-C07-S08-02}
\begin{chapterexercisesolution}{ex:V1-C07-S08-02}
目标$f(x)=x^2$凸，标准约束$c(x)=x-1\le0$仿射，可行域$\mathcal C=(-\infty,1]$非空。取$\widetilde x=0$，
\[
\begin{aligned}
c(\widetilde x)&=0-1=-1<0,\\
\exists\widetilde x:\ c(\widetilde x)<0
&\Longrightarrow \text{Slater条件成立},\\
x^\star&=0,\qquad c(x^\star)=-1<0,\qquad \alpha^\star=0.
\end{aligned}
\]
因此该凸问题强对偶，KKT条件必要且充分。严格可行点只用于资格条件，不必预先等于最优点；本例恰好相同。
\end{chapterexercisesolution}

\Needspace{6\baselineskip}
\begin{exercise}\label{ex:V1-C07-S09-01}
求 $\min_x\tfrac12x^2$，满足$x\ge1$的KKT解与最优值。
\end{exercise}
\phantomsection\label{sol:V1-C07-S09-01}
\begin{chapterexercisesolution}{ex:V1-C07-S09-01}
目标$f(x)=\tfrac12x^2$凸，约束$c(x)=1-x\le0$定义可行域$[1,\infty)$；取$x=2$严格可行，故Slater成立。完整KKT系统为
\[
\begin{aligned}
L(x,\alpha)&=\frac12x^2+\alpha(1-x),\\
\text{驻点：}\quad &x-\alpha=0,\\
\text{原可行：}\quad &1-x\le0,\\
\text{对偶可行：}\quad &\alpha\ge0,\\
\text{互补松弛：}\quad &\alpha(1-x)=0.
\end{aligned}
\]
若$\alpha=0$则驻点给$x=0$，违反原可行性；故约束活动：
\[
x^\star=1,\qquad \alpha^\star=1,\qquad f(x^\star)=\frac12.
\]
凸性与Slater条件使该KKT点为全局最优点。
\end{chapterexercisesolution}

\Needspace{6\baselineskip}
\begin{exercise}\label{ex:V1-C07-S09-02}
在KKT点处若$c_i(x^*)<0$，证明$\alpha_i^*=0$；其逆命题是否成立？
\end{exercise}
\phantomsection\label{sol:V1-C07-S09-02}
\begin{chapterexercisesolution}{ex:V1-C07-S09-02}
在满足KKT的点$(x^\star,\alpha^\star)$，原可行、对偶可行及互补松弛给出
\[
\begin{aligned}
c_i(x^\star)&\le0,\qquad \alpha_i^\star\ge0,\\
\alpha_i^\star c_i(x^\star)&=0,\\
c_i(x^\star)<0&\Longrightarrow \alpha_i^\star=0.
\end{aligned}
\]
逆命题不成立。反例取凸问题$\min_x x^2$，满足$x\ge0$，即$c(x)=-x\le0$。其KKT解为
\[
x^\star=0,\qquad
\alpha^\star=0,\qquad
c(x^\star)=0,
\]
因为驻点$2x-\alpha=0$成立。乘子为零但约束恰好活动。
\end{chapterexercisesolution}

\end{SLChapterExercisePairSection}

\phantomsection\label{struct:V1-C07-CH36}

\Needspace{4\baselineskip}
% KNOWLEDGE-PLAN: KN-V1-C07-CONCEPT-020 L1
\paragraph{读前自检：本章结论与下一步。}凸优化与拉格朗日对偶章末由“目标、可行域、凸性、乘子与对偶”落到“统一约束方向，检查原/对偶可行、驻点和互补松弛”；请把前者产出和后者输入逐项对齐，固定核对前者末项的输出类型能否作为后者首项的输入类型。\par

\begin{SLCompactClosure}
\SLChapterFiveSummary{目标、可行域、凸性、乘子与对偶}{从拉格朗日下确界推出弱对偶，并在资格条件下使用KKT}{统一约束方向，检查原/对偶可行、驻点和互补松弛}{证明候选解全局最优或构造可核验下界}{进入\cref{chap:V1-C08}的数值求解，并保留条件证书}

\phantomsection\label{struct:V1-C07-CH37}
\begin{selfcheckbox}
\begin{enumerate}[leftmargin=2.2em]
  \item 能否把任意方向的不等式统一成$c_i(x)\le0$并正确规定乘子符号？
  \item 能否用线段定义、一阶下界和Hessian三种方式检验适当对象的凸性？
  \item 能否逐行证明弱对偶并说明Slater条件中的严格可行含义？
  \item 能否对一个候选解完整检查五类KKT条件并解释活动约束？未通过时依次回看\cref{sec:V1-C07-S01,sec:V1-C07-S03,sec:V1-C07-S04,sec:V1-C07-S05,sec:V1-C07-S07,sec:V1-C07-S08,sec:V1-C07-S09}。
\end{enumerate}
\end{selfcheckbox}
\end{SLCompactClosure}
